% Internal working draft: synthetic teaching material; no venue selected.
% LaTeX passage; requires amsmath and amssymb in the enclosing document.
\subsection{Retaining fixed linear predictions with null-space updates}
\label{sec:null-space-retention}

How can a linear predictor incorporate an update while retaining its
predictions on a fixed set of old examples? Let $X\in\mathbb{R}^{n\times d}$
contain their feature vectors as rows, and let $w\in\mathbb{R}^d$ denote the
current parameters, so the old prediction vector is $Xw$. An unspecified
learning step supplies a candidate update $u\in\mathbb{R}^d$; its derivation
is outside this passage. The idealized method keeps only the component of
$u$ in the null space $\ker(X)$, the set of parameter directions that leave
these predictions unchanged.

\paragraph{Proposition (exact retention and thresholded retention).}
Fix $X\in\mathbb{R}^{n\times d}$ and let
$N\in\mathbb{R}^{d\times r}$ have orthonormal columns spanning $\ker(X)$,
where $r=\dim\ker(X)$. For every $w,u\in\mathbb{R}^d$, the projected update
satisfies
\begin{equation}
    v=NN^\top u,\qquad w'=w+v,\qquad Xw'=Xw.
\end{equation}
If $r=0$, $N$ has no columns and $v=0$.
For any $\varepsilon>0$, let $P_\varepsilon\in\mathbb{R}^{d\times d}$
be the orthogonal projector onto the right-singular subspace of $X$
with singular values at most $\varepsilon$, including the entire exact
null space. In exact arithmetic, the update $v_\varepsilon=P_\varepsilon u$
satisfies
\begin{equation}
    \|XP_\varepsilon\|_2\leq\varepsilon,
    \qquad
    \|X(w+v_\varepsilon)-Xw\|_2
    \leq\varepsilon\|u\|_2.
\end{equation}
Here $\|\cdot\|_2$ denotes the Euclidean norm for vectors and the spectral
operator norm for matrices.

\paragraph{Why the statements hold.}
Exact retention follows immediately from $XN=0$:
$X(w+NN^\top u)=Xw$. For the thresholded update, the singular value
decomposition shows that $XP_\varepsilon$ retains only directions with
singular values at most $\varepsilon$. Its operator norm therefore has
the stated bound, which gives the prediction-shift bound by multiplication
with $u$.

\paragraph{Relationship to the implementation.}
The implementation uses $P_\varepsilon$ with $\varepsilon>0$, allowing
updates along directions to which the stored predictions are only weakly
sensitive. Such directions can have nonzero singular values, so exact
retention need not hold. The bound limits the total prediction change on
the rows of $X$ relative to the candidate update's size; actual
floating-point numerical error has not been separately bounded.

For an exact-arithmetic illustration, take $X=\operatorname{diag}(1,0.01)$,
$w=(2,3)^\top$, $u=(1,4)^\top$, and $\varepsilon=0.02$.
The exact null space is trivial, so the idealized update is zero.
The thresholded projector keeps the second coordinate, producing
$v_\varepsilon=(0,4)^\top$ and prediction shift $(0,0.04)^\top$.
This is an illustration, not experimental evidence.

Projection may remove a direction needed to learn the new task. Neither
statement guarantees new-task improvement, task accuracy, retention on
other old examples, or generalization to future data. These elementary
linear-algebra properties concern only predictions on the fixed rows of
$X$; they provide no guarantee for nonlinear neural networks.
